\documentclass[a4paper,10pt,5p]{elsarticle}
\usepackage[haranoaji]{luatexja-preset}

\usepackage{amsmath,amssymb}
\usepackage{graphicx}
\usepackage{booktabs}
\usepackage{hyperref}
\usepackage{multirow}      % 行をまたがるセル
\usepackage{siunitx}       % 数字の桁揃え(任意)
\usepackage{xcolor}        % 色指定（周期表のグレー表示）




\graphicspath{{figures2/}}

\journal{Materials Science and Engineering: A}

\begin{document}

\begin{frontmatter}

\title{Effective Atomic Radii in Binary Intermetallic Compounds Using First-Principles Databases}

\author[nims]{Satoshi Minamoto}
\ead{minamoto.satoshi@nims.go.jp}
\author[nims]{Miya Kojima}  
\author[nims]{Akihiro Endo}
% \author[hokudai]{Seiji Miura}



\affiliation[nims]{Materials Data Platform, National Institute for Materials Science (NIMS), Tsukuba, Japan}
% \affiliation[hokudai]{Faculty of Engineering, Hokkaido University, Sapporo, Japan}

\begin{abstract}

B2型(CsCl型)およびL1$_2$型(Cu$_3$Au型)二元金属間化合物における有効原子半径を、Materials Project (MP)およびOpen Quantum Materials Database (OQMD)の格子定数データに対する最適化により決定した。非金属元素およびアルカリ金属元素を含む化合物と格子定数範囲外(2.0--8.0~\AA)の超格子を除外し、MP-B2: 346件、MP-L1$_2$: 723件、OQMD-B2: 2207件、OQMD-L1$_2$: 235件、計3511件を解析した。B2構造にはA--A, B--B, A--B接触条件のmax基準モデル、L1$_2$構造にはA--A, A--B, B--B接触条件のmax基準モデルを適用し、不等式制約付きSLSQP最適化と接触条件の再評価を伴う反復最適化により全元素の有効半径を決定した。Bootstrap法(1000回リサンプリング)による不確定性評価も実施した。

B2とL1$_2$の間で有効半径に系統的な差異が認められた。Re($-5.9$\%)やCa($-5.6$\%)はB2半径がL1$_2$半径より小さく、Ni($+7.9$\%)やPu($+7.4$\%)は逆にB2半径が大きい。遷移金属群では3d: $+1.7$\%、4d: $+0.7$\%、5d: $-0.3$\%と群ごとに異なり、Goldschmidtの一律3\%収縮則とは合致しない。L1$_2$半径はGoldschmidt半径と$R^2=0.858$の相関を示した。非金属・アルカリ金属除外により硬球モデルの適用範囲が明確になった。

得られた構造特異的有効半径は、格子定数予測やサイズミスマッチ評価の基準値として利用できる。
\end{abstract}

\begin{keyword}
atomic radius \sep intermetallic compounds \sep B2 structure \sep L1$_2$ structure \sep first-principles calculation \sep Materials Project \sep OQMD \sep Goldschmidt radius \sep Pauling radius
\end{keyword}

\end{frontmatter}




%%%%%%%%%%%%%%%%%%%%%%%%%%%%%%%%%%%%%
\section{Introduction}

\subsection{Background and Motivation}
原子半径は、格子定数、固溶体形成、相安定性、格子不整合などを支配する基本パラメータであり、Hume--Rothery以来、金属材料学における中心概念として扱われてきた~\citep{humerothery1936}。ただし、原子半径は一意な物理量ではなく、金属半径、共有結合半径、イオン半径など、結合環境に応じた有効量として定義される。金属および金属間化合物に対して最も広く用いられてきた尺度は、Goldschmidt金属半径とPauling金属半径である。

Goldschmidt半径は、純金属の最近接原子間距離を接触剛体球モデルで解釈した経験的尺度であり、配位数12を基準とし、CN=8への変化に対して約3\%の収縮を仮定する。一方、Pauling半径は結合次数と電気陰性度差を考慮したより化学結合論的な尺度である~\citep{pauling1947}。両者は全体として高い相関を示すが、アルカリ金属のような拡散した外殻電子をもつ元素では系統差が見られる。

\begin{figure}[htbp]
    \centering
    \includegraphics[width=0.4\textwidth]{gs_03_pauling_vs_goldschmidt.png}
    \caption{Comparison of Pauling and Goldschmidt metallic radii for elements with available data in both scales. The dashed line represents perfect agreement. Both scales show high correlation ($R^2$=0.993, RMSE=0.026~\AA), but Pauling radii are systematically larger than Goldschmidt radii for alkali metals (Li, Na, K, Rb, Cs).}
    \label{fig:pauling_goldschmidt}
\end{figure}

Slater~\citep{slater1964atomic}は1200以上の結合距離データから半径を導き、平均誤差0.12~\AA を達成した。この半径は共有結合・金属結合・イオン結合のいずれにも適用可能である。

近年、Materials Project (MP)~\citep{jain2013} や Open Quantum Materials Database (OQMD)~\citep{saal2013} などの第一原理データベースが整備され、大規模な格子定数データに基づいて古典的な原子半径尺度を再検討できるようになった。本研究では、B2およびL1$_2$二元金属間化合物の格子定数に対して幾何学モデルを当てはめ、誤差最小化により構造特異的な有効原子半径を決定する。


\subsection{Research Objectives}
B2型(CsCl型, CN=8)とL1$_2$型(Cu$_3$Au型, CN=12)は、NiAl、TiNi、FeAl、Ni$_3$Alなどに代表される重要な秩序相であり、耐熱材料、形状記憶合金、Ni基超合金、高エントロピー合金に広く現れる。両構造は配位数と近接原子配置が大きく異なるため、原子半径の構造依存性を検証するうえで格好の比較対象である。
これを目的として、本研究の目的は以下の4点とする。
\begin{enumerate}
\setlength{\parskip}{0cm}
\setlength{\itemsep}{0cm}
    \item MPおよびOQMDからB2およびL1$_2$二元金属間化合物を収集し、非金属元素除外と格子定数範囲フィルタを適用した解析用データセットを構築する。
    \item B2およびL1$_2$構造に対し、格子定数予測誤差を最小化する構造特異的有効原子半径を最適化により決定する。
    \item 得られた半径をGoldschmidt半径およびPauling半径と比較し、構造依存性と元素群依存性を明らかにする。
    \item MPとOQMDの結果を比較し、決定した半径の頑健性と材料設計への有用性を評価する。
\end{enumerate}


%%%%%%%%%%%%%%%%%%%%%%%%%%%%%%%%%%%%%
\section{Methods}

\subsection{Data Extraction and Curation}

\subsubsection{Materials Project (MP)}

Materials Project (MP)~\citep{jain2013}は、ローレンス・バークレー国立研究所を中心に運営される第一原理計算データベースであり、15万種以上の無機材料の計算特性を収録している。

MPにおける全ての計算は、Vienna Ab initio Simulation Package(VASP)~\citep{kresse1996}で実装された密度汎関数理論(DFT)を用いて行われる。計算では、コア電子の処理にProjector Augmented Wave (PAW)法を採用し、交換相関関数としてPerdew-Burke-Ernzerhof (GGA-PBE)汎化勾配近似を用いる。この一貫した計算手法により、異なる化合物間の有意義な比較が可能となる。

MPにおける主要な計算パラメータには、元素依存性を持つ平面波エネルギーカットオフ値(一般に構成元素の推奨最大カットオフ値の1.3倍)と、原子当たり少なくとも1000個の$k$点密度を持つ$k$点メッシュが含まれる。全構造は、イオン位置と格子定数の両方が、力が0.02 eV/ \AA 以下、応力成分が0.5 kbar以下に収束するまで完全に緩和される。



% The Materials Project ~\citep{jain2013} is a comprehensive first-principles calculation database operated primarily by Lawrence Berkeley National Laboratory, with contributions from numerous academic and industrial partners worldwide. As of the data extraction date, MP contained computed properties for over 150,000 inorganic materials, making it one of the largest and most widely used materials databases.

% All calculations in MP are performed using density functional theory (DFT) as implemented in the Vienna Ab initio Simulation Package (VASP) ~\citep{kresse1996}. The calculations employ the Projector Augmented Wave (PAW) method for treating core electrons and the Generalized Gradient Approximation with the Perdew-Burke-Ernzerhof (GGA-PBE) exchange-correlation functional ~\citep{perdew1996}. This consistent computational methodology across all entries enables meaningful comparisons between different compounds.

% Key computational parameters in MP include plane-wave energy cutoffs that are element-dependent but generally 1.3 times the maximum recommended cutoff for the constituent elements, and k-point meshes with density of at least 1000 k-points per reciprocal atom. All structures are fully relaxed (both ionic positions and cell parameters) until forces converge below 0.02~eV/\AA\ and stress components below 0.5~kbar.

Data were retrieved from the Materials Project database via the Materials Project API (version 2) using the following criteria:
\begin{itemize}
\setlength{\parskip}{0cm}
\setlength{\itemsep}{0cm}
    \item B2 phase: space group Pm$\bar{3}$m (\#221), AB stoichiometry, 2 atoms per cell
    \item L1$_2$ phase: space group Pm$\bar{3}$m (\#221), A$_3$B stoichiometry, 4 atoms per cell
\end{itemize}

以下の品質フィルタを適用した：
\begin{itemize}
\setlength{\parskip}{0cm}
\setlength{\itemsep}{0cm}
    \item \textbf{格子定数範囲フィルタ}(2.0--8.0~\AA)：Rh$_{17}$S$_{15}$($a = 9.95$~\AA)やPd$_{17}$Se$_{15}$($a = 10.70$~\AA)のような超格子構造を除外するため。これらはB2と同じ空間群を持つが、実際には多原子ユニットセルの超格子である。
    \item \textbf{非金属元素除外}：H, He, C, N, O, F, Ne, P, S, Cl, Ar, Se, Br, Kr, Te, I, Xe, At, Rnを含む化合物を除外。これらの元素は共有結合やイオン結合を形成する傾向が強く、金属間化合物を対象とした硬球モデルの適用範囲外となる。例として、IrN($a_{\mathrm{DFT}} = 4.10$~\AA, $a_{\mathrm{calc}} = 2.87$~\AA, 誤差$-30$\%)やTiO($a_{\mathrm{DFT}} = 4.09$~\AA, $a_{\mathrm{calc}} = 3.11$~\AA, 誤差$-24$\%)のような窒化物・酸化物は、DFT格子定数に対して硬球モデルが著しい過小評価を示す。
\end{itemize}

% Instead of the conventional energy\_above\_hull threshold filter commonly used in high-throughput screening studies, this study adopted a different filtering strategy. The energy\_above\_hull filter was not applied for two reasons. First, effective radius optimization is a statistical method that simultaneously fits hundreds to thousands of compounds, diluting the influence of individual unstable compounds on the overall result. Second, application of the stability filter can systematically exclude certain element combinations (e.g., alkali metal--non-metal B2 compounds), introducing bias in the elemental coverage of the dataset.

% Instead, the following quality filters were applied: (1) a lattice constant range filter (2.0--8.0~\AA) to exclude superstructures such as Rh$_{17}$S$_{15}$ ($a = 9.95$~\AA) and Pd$_{17}$Se$_{15}$ ($a = 10.70$~\AA), which share the same space group as B2 but are actually multi-atom superstructures; and (2) exclusion of compounds containing non-metallic elements (H, He, C, N, O, F, Ne, P, S, Cl, Ar, Se, Br, Kr, Te, I, Xe, At, Rn), as these elements tend to form covalent or ionic bonds that fall outside the applicability of the hard-sphere model for intermetallic compounds.

抽出結果、MPからの取得件数はB2系490件、L1$_2$系932件であった。品質フィルタ(格子定数範囲および非金属元素除外)を適用後、分析対象としてB2系357件(68元素)、L1$_2$系848件(72元素)が残った。非金属フィルタの適用により、特にMP-B2ケースではNa、Kなどのアルカリ金属のB2化合物がほぼ全て除外された。これは、MPにおけるNaおよびKのB2構造化合物がNaCl、KBr等の非金属パートナーとの化合物に集中しているためである。このことは、これらの化合物が実質的にイオン性であり、硬球モデルの適用対象外であることを示唆している。抽出結果は表~\ref{tab:combined_data}にまとめられている。

% The extraction yielded 490 B2 entries and 932 L1$_2$ entries from MP before filtering. After applying the quality filters (lattice constant range and non-metallic element exclusion), 357 B2 compounds (68 elements) and 848 L1$_2$ compounds (72 elements) remained for analysis. The non-metallic filter notably excluded nearly all B2 compounds of Na and K alkali metals, because their B2 structures in MP are concentrated among compounds with non-metallic partners such as NaCl and KBr, indicating that these are essentially ionic compounds outside the applicability of the hard-sphere model.

\subsubsection{Open Quantum Materials Database (OQMD)}
OQMD~\citep{saal2013}は、ノースウェスタン大学を中心に開発・維持されている主要な第一原理計算データベースである。OQMDは熱力学的性質と相安定性に重点を置いて設計され、多数の化合物について密度汎関数理論(DFT)で計算された生成エネルギーと構造パラメータを含む。MPと同様に、OQMDの計算はVASPを用いてPAWポテンシャルとGGA-PBE関数で実行されるため、両データベース間の広範な互換性が確保されている。ただし、k点密度や収束基準など、MPとOQMDでは計算パラメータに差異があり、計算特性にわずかな系統的差異が生じる可能性がある。

OQMDからのデータ抽出はREST APIを通じて行い、B2構造にはプロトタイプ``B2\_CsCl''、L1$_2$構造にはプロトタイプ``AuCu3''を指定した。フィルタリング前の取得件数はB2系3469件、L1$_2$系273件であった。MPと同じ品質フィルタ(格子定数範囲2.0--8.0~\AA、非金属元素除外)を適用後、分析対象としてB2系2539件(72元素)、L1$_2$系238件(58元素)が保持された。OQMDデータセットはMPよりもB2化合物を大幅に多く含む(2539対357)一方で、MPはL1$_2$化合物をより多く含む(848対238)点が特徴的であり、これはデータベース構築における重点の違いを反映している。

% The Open Quantum Materials Database (OQMD) ~\citep{saal2013} is another major first-principles calculation database, developed and maintained primarily at Northwestern University. Data extraction from OQMD was performed via the REST API, specifying prototype ``B2\_CsCl'' for B2 structures and ``AuCu3'' for L1$_2$ structures. The extraction yielded 3469 B2 entries and 273 L1$_2$ entries before filtering. After applying the same quality filters as MP (lattice constant range 2.0--8.0~\AA\ and non-metallic element exclusion), 2539 B2 compounds (72 elements) and 238 L1$_2$ compounds (58 elements) were retained for analysis.


\begin{table*}[t]
  \small
  \centering
  \caption{Summary of data extracted from the Materials Project and OQMD databases for B2 and L1$_{2}$ structures. The ``Retrieved'' column shows the initial number of entries matching the structure type query. ``After filter'' indicates the entries remaining after applying the lattice constant range filter (2.0--8.0~\AA) and non-metallic element exclusion. ``Elements'' represents the number of unique elements in the filtered dataset, and ``Representative compounds'' lists examples of compounds for each group.}
  \label{tab:combined_data}
  \scalebox{0.93}{%
\scalebox{0.9}{
  \begin{tabular}{llcccc}
    \toprule
    \multirow{2}{*}{Database} & \multirow{2}{*}{Structure} &
    \multicolumn{2}{c}{Number of entries} &
    \multicolumn{2}{c}{Composition} \\
    \cmidrule(lr){3-4} \cmidrule(lr){5-6}
    & & Retrieved & After filter & Elements & Representative compounds \\
    \midrule
    \multirow{2}{*}{Materials Project}%
    & B2   & 490 & 357 & 68 & NiAl, FeAl, TiNi, CoAl \\
    & L1\(_2\) & 932 & 848 & 72 & Ni\(_3\)Al, Ni\(_3\)Ti, Co\(_3\)Ti, Pt\(_3\)Al \\
    \midrule
    \multirow{2}{*}{OQMD}%
    & B2   & 3469 & 2539 & 72 & CrAl, MoAl, WAl, NbAl \\
    & L1\(_2\) & 273 & 238 & 58 & Cr\(_3\)Al, Mo\(_3\)Al, W\(_3\)Al \\
    \bottomrule
  \end{tabular}}}
\end{table*}



\subsubsection{Dataset Comparison}

MPとOQMDのデータセットサイズの比較を図~\ref{fig:dataset_size}に示す。両データベースの補完性、つまりOQMDはB2データが豊富(2539対357)である一方、MPはL1$_2$エントリが多い(848対238)ことにより、結果の交差検証が可能となり、異なる計算プロトコル間で決定された半径の頑健性を評価できる。総計3982件のデータセットを用いて解析を行った。

% A comparison of dataset sizes between MP and OQMD is shown in Fig.~\ref{fig:dataset_size}. The complementary nature of the two databases---OQMD being richer in B2 data (2539 vs. 357) while MP has more L1$_2$ entries (848 vs. 238)---provides an opportunity to cross-validate results and assess the robustness of the determined radii across different computational protocols. A total of 3982 compounds were analyzed.

\begin{figure}[h]
    \centering
    \includegraphics[width=0.4\textwidth]{00a_dataset_size.png}
    \caption{Comparison of dataset sizes between Materials Project (MP) and OQMD for B2 and L1$_2$ structures. OQMD contains substantially more B2 compounds (2539 vs. 357), while MP has more L1$_2$ compounds (848 vs. 238). This complementary coverage enables cross-validation of results between databases.}
    \label{fig:dataset_size}
\end{figure}

\subsection{Geometric Models for Lattice Constants}

格子定数と原子半径の関係は、結晶構造と原子配列に関する仮定に依存する。ここで検討するB2およびL1$_2$構造については、特定の結晶学方向に沿って原子が接触する剛体球充填に基づく幾何学的モデルを採用する。

% The relationship between lattice constants and atomic radii depends on the crystal structure and the assumptions made about atomic packing. For the B2 and L1$_2$ structures considered here, we adopt geometric models based on hard-sphere packing with atoms in contact along specific crystallographic directions.

\subsubsection{B2 Structure Model}

B2(CsCl型)構造は、互いに重複する二つの単純立方格子サブ格子から成る。A原子は立方体の頂点(0,0,0)に位置し、B原子は体心(1/2,1/2,1/2)に位置する。両サブ格子とも同種原子間の最近接距離は格子定数$a$に等しい。A--B間の最近接距離は体対角線の半分$a\sqrt{3}/2$である。したがって、B2構造では三つの接触条件が考えられる：

\begin{enumerate}
\setlength{\parskip}{0cm}
\setlength{\itemsep}{0cm}
    \item A--A接触（立方体の辺に沿った同種原子間接触）: $a_1 = 2r_A$
    \item B--B接触（隣接単位格子間のB原子間接触）: $a_2 = 2r_B$
    \item A--B接触（体対角線に沿った異種原子間接触）: $a_3 = \frac{2(r_A + r_B)}{\sqrt{3}}$
\end{enumerate}

剛体球モデルでは原子の重なりは許されないため、格子定数は三つの接触条件のうち最大の値によって決定される：

\begin{equation}
    a_{B2} = \max\!\left(2r_A,\; 2r_B,\; \frac{2(r_A + r_B)}{\sqrt{3}}\right)
    \label{eq:B2}
\end{equation}

ここで、$r_A$および$r_B$はそれぞれA種とB種の原子半径である。$r_A$と$r_B$のサイズ差が大きい場合には同種原子間接触(A--AまたはB--B)が支配的となり、両者のサイズが近い場合にはA--B接触条件が支配的となる。このモデルはA/Bのラベリングに依存せず、$\max(2r_A, 2r_B) = 2\max(r_A, r_B)$と等価である。

\subsubsection{L1$_2$ Structure Model}

L1$_2$(Cu$_3$Au型)構造は、面中心にA原子(多数種)、立方体の角にB原子(少数種)を配置した、面心立方構造の秩序ある派生形である。配位環境はB2よりも複雑である：A原子は8つのA近接原子と4つのB近接原子(全て距離$\frac{a}{\sqrt{2}}$)を持つ一方、B原子は同じ距離に12のA近接原子を持つ。

L1$_2$では、A原子とB原子の相対的な大きさに依存して複数の接触状態が取り得る：

% The L1$_2$ (Cu$_3$Au-type) structure is an ordered derivative of the FCC structure, with A atoms (majority species) at face centers and B atoms (minority species) at cube corners. The coordination environment is more complex than B2: A atoms have 8 A neighbors and 4 B neighbors (all at distance $\frac{a}{\sqrt{2}}$), while B atoms have 12 A neighbors at the same distance.

% For L1$_2$, multiple contact conditions are possible depending on the relative sizes of A and B atoms:

\begin{enumerate}
\setlength{\parskip}{0cm} % 段落間
\setlength{\itemsep}{0cm} % 項目間
    \item A-A contact along face diagonals: $a_1 = 2\sqrt{2}\, r_A$
    \item A-B contact along face diagonals: $a_2 = \sqrt{2}(r_A + r_B)$
    \item B-B contact along cube edges: $a_3 = 2r_B$
\end{enumerate}

実際の結晶構造では、原子サイズの組み合わせに応じていずれかの接触条件が格子定数を支配する。$r_A$が$r_B$より十分大きい場合はA-A接触条件が支配的であり、$r_B$が非常に大きい場合はB-B接触条件が支配的となる。剛体球モデルでは、原子の重なりは許されないため、格子定数は三つの接触条件のうち最大のものによって決定される：

% In real crystal structures, one of the contact conditions dominates the lattice parameter depending on the combination of atomic sizes. When $r_A$ is sufficiently larger than $r_B$, the A-A contact condition dominates, while the B-B contact condition becomes dominant when $r_B$ is very large. In the hard-sphere model, atomic overlap is not allowed, so the lattice parameter is determined by the maximum of the three contact conditions:

\begin{equation}
    a_{L1_2} = \max(a_1,\, a_2,\, a_3) = \max\!\left(2\sqrt{2}\,r_A,\; \sqrt{2}(r_A + r_B),\; 2r_B\right)
    \label{eq:L12}
\end{equation}

ここで$r_A$は多数種(A$_3$B中のA)の半径、$r_B$は少数種の半径である。この最大値モデルは、剛体球が互いに重なることなく結晶格子内に収まるために必要な最小の格子定数を与えるという物理的描像に基づいている。A原子のサイズが支配的な場合はA-A接触($a_1$)、A-B間のサイズ和が大きい場合はA-B接触($a_2$)、B原子が特に大きい場合はB-B接触($a_3$)が格子定数を決定する。

% where $r_A$ is the radius of the majority species (A in A$_3$B) and $r_B$ is the radius of the minority species. This maximum model is based on the physical picture that the lattice parameter must be large enough to accommodate hard spheres without overlap. When A atoms are dominant in size, A-A contact ($a_1$) determines the lattice parameter; when the A-B size sum is large, A-B contact ($a_2$) dominates; when B atoms are particularly large, B-B contact ($a_3$) determines the lattice parameter.

B2構造とL1$_2$構造を示す模式図は
% 図~\ref{fig:structure}に示され、3D透視図は図
~\ref{fig:structure_3d}に示されている。

% Schematic diagrams illustrating the B2 and L1$_2$ structures are shown in Fig.~\ref{fig:structure}, and 3D perspective views are provided in Fig.~\ref{fig:structure_3d}.

\begin{figure*}[htbp]
    \centering
    \includegraphics[width=0.5\textwidth]{22_structure_3d_multiview.png}
    \caption{Schematic diagrams of B2 (CsCl-type) and L1$_2$ (Cu$_3$Au-type) crystal structures. In B2, A atoms (blue) occupy cube corners and B atoms (red) occupy body centers, with each atom having 8 nearest neighbors of the opposite type. In L1$_2$, A atoms (blue) occupy face-centered positions and B atoms (red) occupy cube corners, with A atoms having 12 nearest neighbors (8 A + 4 B) and B atoms having 12 A neighbors. Three-dimensional perspective views of B2 and L1$_2$ unit cells, showing the spatial arrangement of atoms and the nearest-neighbor bonding environment.}
    \label{fig:structure_3d}
\end{figure*}

\subsection{Radius Optimization Algorithm}

格子定数と原子半径の関係を示す幾何学的モデルを考慮すると、有効半径を決定する作業は最適化問題となる。すなわち、データセット内の全化合物において観測値(DFT計算値)と予測値の格子定数との差異を最小化する半径の集合$\{r_i\}$を求める問題である。

これを非線形最小二乗問題として定式化する：

% Given the geometric models relating lattice constants to atomic radii, the task of determining effective radii becomes an optimization problem: find the set of radii $\{r_i\}$ that minimizes the discrepancy between observed (DFT-calculated) and predicted lattice constants across all compounds in the dataset.

% We formulate this as a nonlinear least-squares problem:

\begin{equation}
    \min_{\{r_i\}} \sum_{j=1}^{N} \left( a_j^{obs} - a_j^{pred}(\{r_i\}) \right)^2
    \label{eq:optimization}
\end{equation}

ここで、$N$は化合物の数、$a_j^{obs}$は化合物$j$の観測された格子定数、$a_j^{pred}$は現在の半径セットを用いて幾何学的モデルから計算された予測格子定数である。

B2のmax基準モデル（式~\ref{eq:B2}）およびL1$_2$のmax基準モデル（式~\ref{eq:L12}）では、$\max$関数の非微分点が最適化ソルバーに問題を生じさせる。これを回避するため、接触条件の再評価を伴う反復最適化スキームを採用した。

各反復では、(1) 現在の半径から各化合物の支配的接触条件を決定し、(2) その接触条件を固定した最適化部分問題を解き、(3) 新しい半径で接触条件を再評価する。接触条件の切り替わりがなくなった時点で収束とみなす。この反復スキームにより、各化合物の接触条件がその最適化された半径と自己無撞着になることが保証される。

最適化にはSciPyライブラリのSequential Least Squares Programming (SLSQP)アルゴリズムを使用した。全ての半径に対して物理的制約($0.5 \leq r_i \leq 3.5$~\AA)を課し、さらに非支配的接触条件が観測格子定数以下となる不等式制約を導入した。これにより、硬球が重ならないという物理的条件が保証される。初期値にはPaulingの金属半径を用いた。最適化はB2構造とL1$_2$構造で別々に実行され、各元素ごとに構造固有の有効半径が得られた。

半径の不確定性はBootstrap法により評価した。各ケースについて1000回のリサンプリングを行い、各リサンプルに対して同一の反復最適化を実行して半径分布を得た。各元素の半径の標準偏差および95\%信頼区間を算出した。


% where $N$ is the number of compounds, $a_j^{obs}$ is the observed lattice constant for compound $j$, and $a_j^{pred}$ is the predicted lattice constant calculated from the geometric model using the current set of radii.

% The optimization was performed using the Trust Region Reflective (TRF) algorithm as implemented in the SciPy library. TRF is well-suited for bound-constrained optimization problems and provides robust convergence even for ill-conditioned systems. We imposed physical constraints requiring all radii to be positive ($r_i > 0$) and bounded the search space to reasonable values (0.5--3.5~\AA) based on known atomic sizes.

% Initial radius values were taken from Pauling metallic radii where available, with interpolated or estimated values for elements not in Pauling's original tabulation. The optimization was run separately for B2 and L1$_2$ structures, yielding structure-specific effective radii for each element.

% Convergence was assessed by monitoring the reduction in the objective function (sum of squared residuals) and the change in radius values between iterations. The optimization was considered converged when the relative change in the objective function fell below $10^{-8}$ and the maximum change in any radius was less than $10^{-6}$~\AA.

\subsection{Reference Radii for Comparison}

The optimized effective radii were compared with two classical radius scales:
\begin{itemize}
\setlength{\parskip}{0cm} % 段落間
\setlength{\itemsep}{0cm} % 項目間
    \item \textbf{Goldschmidt金属半径} ~\citep{goldschmidt1927}: 純金属結晶構造における原子間距離から導かれた経験的半径。配位数12を基準とする。値は標準的な集計データから引用した。
   \item \textbf{Pauling金属半径} ~\citep{pauling1947}: 価電子結合理論と結合次数-結合距離関係に基づく理論半径。値はPaulingの1947年原著及びその後の更新版から引用した。
\end{itemize}


% \begin{itemize}
%    \item \textbf{Goldschmidt metallic radii} ~\citep{goldschmidt1927}: Empirical radii derived from interatomic distances in pure metal crystal structures, referenced to coordination number 12. Values were taken from standard compilations.
%    \item \textbf{Pauling metallic radii} ~\citep{pauling1947}: Theoretical radii based on valence bond theory and the bond order--bond distance relationship. Values were taken from Pauling's original 1947 publication and subsequent updates.
% \end{itemize}


Goldschmidt半径とPauling半径の両方が利用可能な元素については、直接比較を行った。相関係数($R^2$)と二乗平均平方根誤差(RMSE)を算出し、最適化半径と参照半径の一致度を定量化した。

% For elements where both Goldschmidt and Pauling radii are available, direct comparisons were made. The correlation coefficient ($R^2$) and root-mean-square error (RMSE) were calculated to quantify the agreement between optimized and reference radii.

\subsection{Statistical Analysis}

Model performance was evaluated using standard regression metrics:
\begin{itemize}
\setlength{\parskip}{0cm} % 段落間
\setlength{\itemsep}{0cm} % 項目間
    \item Root-mean-square error (RMSE): \\$\sqrt{\frac{1}{N}\sum_j(a_j^{obs} - a_j^{pred})^2}$
    \item Coefficient of determination ($R^2$): \\$1 - \frac{\sum_j(a_j^{obs} - a_j^{pred})^2}{\sum_j(a_j^{obs} - \bar{a})^2}$
\end{itemize}

ここで、$\bar{a}$は観測された平均格子定数である。これらの指標は、各データベース(MP、OQMD)および構造タイプ(B2、L1$_2$)ごとに個別に計算された。

元素群分析は、元素をカテゴリー(アルカリ金属、アルカリ土類金属、3d/4d/5d遷移金属、ランタノイド、アクチノイド、遷移金属後元素)に分類し、各群内で平均半径、平均B2-L1$_2$差、標準偏差などの群平均特性を算出することで実施された。

% where $\bar{a}$ is the mean observed lattice constant. These metrics were calculated separately for each database (MP, OQMD) and structure type (B2, L1$_2$).

% Element-group analysis was performed by classifying elements into categories (alkali metals, alkaline earth metals, 3d/4d/5d transition metals, lanthanides, actinides, and post-transition metals) and calculating group-averaged properties such as mean radius, mean B2-L1$_2$ difference, and standard deviation within each group.

%%%%%%%%%%%%%%%%%%%%%%%%%%%%%%%%%%%%%
\section{Results}

\subsection{Model Performance}

最適化された有効半径を用いた格子定数予測の精度を表~\ref{tab:performance}にまとめた。B2構造では、MP-B2のRMSEが0.054~\AA と最も小さく、平均相対誤差は1.03\%であった。OQMD-B2はデータ量が6倍(2207対346)と大幅に多いにもかかわらず、RMSE = 0.095~\AA、平均相対誤差1.90\%と良好な精度を維持している。

L1$_2$構造については、OQMD-L1$_2$がRMSE = 0.066~\AA、平均相対誤差1.13\%と4ケース中最高の精度を示した。MP-L1$_2$はRMSE = 0.116~\AA とやや大きいが、データ量が723件と最大であり、化合物の多様性がより大きいことを反映している可能性がある。$R^2$値は0.935--0.975の範囲にあり、幾何学的モデルが格子定数の変動の大部分を捉えていることを示している。

% The accuracy of lattice constant prediction using the optimized effective radii is summarized in Table~\ref{tab:performance}. For B2 structures, MP-B2 achieved the smallest RMSE of 0.052~\AA\ with a mean relative error of 0.99\%. OQMD-B2, despite having seven times more data (2539 vs. 357), maintained good accuracy with RMSE = 0.089~\AA\ and mean relative error of 1.74\%.

% For L1$_2$ structures, OQMD-L1$_2$ showed the highest accuracy among all four cases with RMSE = 0.066~\AA\ and mean relative error of 1.13\%. MP-L1$_2$ had a somewhat larger RMSE of 0.157~\AA, but its dataset of 848 compounds is the largest, and the greater diversity of compounds may contribute to the larger error. $R^2$ values remained above 0.94 in all cases, indicating that the geometric model captures the majority of variance in lattice constants.

\begin{table*}[htbp]
\small
    \centering
    \caption{Lattice constant prediction accuracy using optimized effective radii. RMSE is the root-mean-square error between observed and predicted lattice constants, MAE is the mean absolute error, and $R^2$ is the coefficient of determination. The mean and maximum relative errors are also shown.}
    \label{tab:performance}
\scalebox{0.85}{
    \begin{tabular}{llccccc}
        \toprule
        Database & Structure & RMSE (\AA) & MAE (\AA) & $R^2$ & Mean rel.\ err.\ (\%) & Max rel.\ err.\ (\%) \\
        \midrule
        MP & B2 & 0.054 & 0.036 & 0.975 & 1.03 & 9.88 \\
        MP & L1$_2$ & 0.116 & 0.093 & 0.950 & 2.08 & 9.52 \\
        OQMD & B2 & 0.095 & 0.067 & 0.935 & 1.90 & 19.20 \\
        OQMD & L1$_2$ & 0.066 & 0.047 & 0.973 & 1.13 & 6.88 \\
        \bottomrule
    \end{tabular}
    }
\end{table*}

図~\ref{fig:performance}は、データベースと構造タイプごとの予測精度をグラフで比較したものである。MP-B2が最小RMSE(0.054~\AA)を示す一方、OQMD-L1$_2$が最小の平均相対誤差(1.13\%)を示しており、データソースと構造タイプの組み合わせによって精度特性が異なることが分かる。

% A graphical comparison of prediction accuracy across databases and structure types is shown in Fig.~\ref{fig:performance}. MP-B2 shows the smallest RMSE while OQMD-L1$_2$ shows the smallest mean relative error, indicating that accuracy characteristics differ depending on the combination of data source and structure type.

\begin{figure}[htbp]
    \centering
    \includegraphics[width=0.45\textwidth]{12_model_performance.png}
    \caption{Comparison of lattice constant prediction accuracy (RMSE) across databases and structure types. MP-B2 achieves the highest accuracy (RMSE = 0.054~\AA).}
    \label{fig:performance}
\end{figure}

\subsection{Bootstrap Uncertainty Estimation}
Bootstrap法(1000回リサンプリング)により各元素の有効半径の不確定性を評価した結果を表~\ref{tab:bootstrap}にまとめた。全4ケースで1000回中1000回の最適化が成功した。半径の標準偏差の中央値はOQMD-B2で最小(0.007~\AA)、OQMD-L1$_2$で最大(0.021~\AA)であった。95\%信頼区間の平均幅はOQMD-B2: 0.036~\AA、MP-B2: 0.065~\AA、MP-L1$_2$: 0.070~\AA、OQMD-L1$_2$: 0.094~\AA であった。

不確定性が大きい元素は、各ケースで異なる傾向を示した。MP-B2ではFe (std = 0.052~\AA)、MP-L1$_2$ではBa (0.040~\AA)、OQMD-B2ではB (0.039~\AA)、OQMD-L1$_2$ではFe (0.060~\AA)が最大の不確定性を示した。OQMD-B2のstdが最小なのはデータ量(2207件)が多いことを反映している。

\begin{table}[htbp]
\small
    \centering
    \caption{Summary of Bootstrap uncertainty estimation (1000 resamples per case). Median and maximum standard deviation of effective radii, and the average width of 95\% confidence intervals are shown.}
    \label{tab:bootstrap}
    \begin{tabular}{lcccc}
        \toprule
        Case & Med.\ std (\AA) & Max.\ std (\AA) & Avg.\ 95\% CI (\AA) & Max.\ std element \\
        \midrule
        MP-B2 & 0.014 & 0.052 & 0.065 & Fe \\
        MP-L1$_2$ & 0.017 & 0.040 & 0.070 & Ba \\
        OQMD-B2 & 0.007 & 0.039 & 0.036 & B \\
        OQMD-L1$_2$ & 0.021 & 0.060 & 0.094 & Fe \\
        \bottomrule
    \end{tabular}
\end{table}

\subsection{Parity Plots}
観測値と予測値のパリティプロットは、B2では対角線近傍に点が密集し、系統的偏差が小さいことを示している。一方、L1$_2$では全体として良好な一致を保つものの、B2よりも散らばりがやや大きい。これはL1$_2$で複数の接触条件が競合し得ること、および剛体球近似で捉えられない電子的寄与を反映している可能性がある。

\begin{figure*}[t]
  \centering
  \begin{minipage}[t]{0.48\textwidth}
    \centering
    \includegraphics[width=0.9\linewidth]{03_B2_parity_lattice.png}
    \caption{Parity plot for MP B2 structures showing observed vs. predicted lattice constants. The dashed line represents perfect agreement. The tight clustering of points around the diagonal demonstrates the high accuracy of the geometric model with optimized radii (RMSE = 0.054~\AA, $R^2 = 0.975$).}
    \label{fig:parity_B2}
  \end{minipage}
  \hfill
  \begin{minipage}[t]{0.48\textwidth}
    \centering
    \includegraphics[width=0.9\linewidth]{04_L12_parity_lattice.png}
    \caption{Parity plot for MP L1$_2$ structures showing observed vs. predicted lattice constants. The agreement is good overall (RMSE = 0.116~\AA, $R^2 = 0.950$), though with somewhat greater scatter than for B2 structures.}
    \label{fig:parity_L12}
  \end{minipage}
\end{figure*}


\subsection{Comparison with Classical Radii}

\subsubsection{Comparison with Goldschmidt Radii}

最適化されたL1$_2$半径は、図~\ref{fig:L12_goldschmidt}に示すように、Goldschmidt半径と相関を示す($R^2 = 0.858$、RMSE = 0.083~\AA)。L1$_2$は最密充填誘導体構造であり主成分原子は12個の最近接原子を持つため、配位数12の純金属から導かれたGoldschmidtの経験値との対応は物理的に合理的である。RMSE = 0.083~\AA の偏差は、純金属と二元金属間化合物における結合環境の違いを反映している。

% The optimized L1$_2$ radii show good correlation with Goldschmidt radii ($R^2 = 0.946$, RMSE = 0.111~\AA), as shown in Fig.~\ref{fig:L12_goldschmidt}.

\begin{figure}[htbp]
    \centering
    \includegraphics[width=0.4\textwidth]{gs_02_L12_vs_goldschmidt.png}
    \caption{Comparison of optimized L1$_2$ radii (MP) with Goldschmidt metallic radii ($R^2 = 0.858$, RMSE = 0.083~\AA, 60 common elements).}
    \label{fig:L12_goldschmidt}
\end{figure}

\subsubsection{Comparison with Pauling Radii}


B2半径とPauling半径の比較は、図~\ref{fig:B2_pauling}に示すように、相関を示す(全62元素: $R^2 = 0.835$)。max基準モデルにおいて同種原子間接触(A--AまたはB--B)が支配的な化合物では、非支配的な原子種の半径は格子定数に直接影響しないため、その最適化に対する制約が弱まることを反映している。

% The comparison of B2 radii with Pauling radii shows good correlation (all 66 elements: $R^2 = 0.925$; excl.~alkali metals (Li, Rb, Cs) 63 elements: $R^2 = 0.913$), as shown in Fig.~\ref{fig:B2_pauling}.

\begin{figure}[htbp]
    \centering
    \includegraphics[width=0.4\textwidth]{05_B2_vs_pauling.png}
    \caption{Comparison of optimized B2 radii (MP) with Pauling metallic radii (62 common elements: $R^2 = 0.835$).}
    \label{fig:B2_pauling}
\end{figure}

我々は、Materials ProjectおよびOQMD第一原理データベースの格子定数データに対して最適化を行うことで、B2型およびL1$_2$型二元金属間化合物の有効原子半径を決定した。
得られた原子半径をB2構造、L1$_2$構造ごとにTable~\ref{tab:radii_b2}、Table~\ref{tab:radii_l12}に示す。各表では上段にMPの値、下段にOQMDの値を記載している。
  \begin{table*}[htbp]
\centering
\scriptsize
\caption{Optimized effective atomic radii (\AA) for B2 structure.
Upper values: MP (357 compounds, 68 elements),
lower values: OQMD (2539 compounds, 72 elements).
Non-metallic elements (H, He, C, N, O, F, Ne, P, S, Cl, Ar, Se, Br, Kr, Te, I, Xe, At, Rn) were excluded from the dataset prior to optimization; only elements appearing in the filtered compound set have numerical values.}
\label{tab:radii_b2}
\setlength{\tabcolsep}{1.5pt}
\begin{tabular}{cccccccccccccccccc}
\toprule
\makebox[0.9cm]{\textbf{H}} &  &  &  &  &  &  &  &  &  &  &  &  &  &  &  &  & \makebox[0.9cm]{\textbf{He}} \\
\midrule
\makebox[0.9cm]{\begin{tabular}{@{}c@{}}\textbf{Li}\\1.32\\1.37\end{tabular}} & \makebox[0.9cm]{\begin{tabular}{@{}c@{}}\textbf{Be}\\1.10\\1.19\end{tabular}} &  &  &  &  &  &  &  &  &  &  & \makebox[0.9cm]{\begin{tabular}{@{}c@{}}\textbf{B}\\1.09\\1.07\end{tabular}} & \makebox[0.9cm]{\textbf{C}} & \makebox[0.9cm]{\textbf{N}} & \makebox[0.9cm]{\textbf{O}} & \makebox[0.9cm]{\textbf{F}} & \makebox[0.9cm]{\textbf{Ne}} \\
\makebox[0.9cm]{\begin{tabular}{@{}c@{}}\textbf{Na}\\--\\1.70\end{tabular}} & \makebox[0.9cm]{\begin{tabular}{@{}c@{}}\textbf{Mg}\\1.54\\1.54\end{tabular}} &  &  &  &  &  &  &  &  &  &  & \makebox[0.9cm]{\begin{tabular}{@{}c@{}}\textbf{Al}\\1.36\\1.41\end{tabular}} & \makebox[0.9cm]{\begin{tabular}{@{}c@{}}\textbf{Si}\\1.30\\1.33\end{tabular}} & \makebox[0.9cm]{\textbf{P}} & \makebox[0.9cm]{\textbf{S}} & \makebox[0.9cm]{\textbf{Cl}} & \makebox[0.9cm]{\textbf{Ar}} \\
\makebox[0.9cm]{\begin{tabular}{@{}c@{}}\textbf{K}\\--\\2.07\end{tabular}} & \makebox[0.9cm]{\begin{tabular}{@{}c@{}}\textbf{Ca}\\1.76\\1.83\end{tabular}} & \makebox[0.9cm]{\begin{tabular}{@{}c@{}}\textbf{Sc}\\1.54\\1.55\end{tabular}} & \makebox[0.9cm]{\begin{tabular}{@{}c@{}}\textbf{Ti}\\1.41\\1.38\end{tabular}} & \makebox[0.9cm]{\begin{tabular}{@{}c@{}}\textbf{V}\\1.31\\1.32\end{tabular}} & \makebox[0.9cm]{\begin{tabular}{@{}c@{}}\textbf{Cr}\\1.32\\1.33\end{tabular}} & \makebox[0.9cm]{\begin{tabular}{@{}c@{}}\textbf{Mn}\\1.31\\1.28\end{tabular}} & \makebox[0.9cm]{\begin{tabular}{@{}c@{}}\textbf{Fe}\\1.22\\1.21\end{tabular}} & \makebox[0.9cm]{\begin{tabular}{@{}c@{}}\textbf{Co}\\1.17\\1.19\end{tabular}} & \makebox[0.9cm]{\begin{tabular}{@{}c@{}}\textbf{Ni}\\1.18\\1.20\end{tabular}} & \makebox[0.9cm]{\begin{tabular}{@{}c@{}}\textbf{Cu}\\1.26\\1.28\end{tabular}} & \makebox[0.9cm]{\begin{tabular}{@{}c@{}}\textbf{Zn}\\1.34\\1.37\end{tabular}} & \makebox[0.9cm]{\begin{tabular}{@{}c@{}}\textbf{Ga}\\1.36\\1.41\end{tabular}} & \makebox[0.9cm]{\begin{tabular}{@{}c@{}}\textbf{Ge}\\--\\1.41\end{tabular}} & \makebox[0.9cm]{\begin{tabular}{@{}c@{}}\textbf{As}\\1.41\\1.41\end{tabular}} & \makebox[0.9cm]{\textbf{Se}} & \makebox[0.9cm]{\textbf{Br}} & \makebox[0.9cm]{\textbf{Kr}} \\
\makebox[0.9cm]{\begin{tabular}{@{}c@{}}\textbf{Rb}\\2.11\\2.21\end{tabular}} & \makebox[0.9cm]{\begin{tabular}{@{}c@{}}\textbf{Sr}\\1.94\\1.99\end{tabular}} & \makebox[0.9cm]{\begin{tabular}{@{}c@{}}\textbf{Y}\\1.71\\1.73\end{tabular}} & \makebox[0.9cm]{\begin{tabular}{@{}c@{}}\textbf{Zr}\\1.57\\1.53\end{tabular}} & \makebox[0.9cm]{\begin{tabular}{@{}c@{}}\textbf{Nb}\\1.51\\1.41\end{tabular}} & \makebox[0.9cm]{\begin{tabular}{@{}c@{}}\textbf{Mo}\\1.43\\1.33\end{tabular}} & \makebox[0.9cm]{\begin{tabular}{@{}c@{}}\textbf{Tc}\\1.27\\1.29\end{tabular}} & \makebox[0.9cm]{\begin{tabular}{@{}c@{}}\textbf{Ru}\\1.24\\1.27\end{tabular}} & \makebox[0.9cm]{\begin{tabular}{@{}c@{}}\textbf{Rh}\\1.26\\1.28\end{tabular}} & \makebox[0.9cm]{\begin{tabular}{@{}c@{}}\textbf{Pd}\\1.31\\1.34\end{tabular}} & \makebox[0.9cm]{\begin{tabular}{@{}c@{}}\textbf{Ag}\\1.42\\1.45\end{tabular}} & \makebox[0.9cm]{\begin{tabular}{@{}c@{}}\textbf{Cd}\\1.52\\1.55\end{tabular}} & \makebox[0.9cm]{\begin{tabular}{@{}c@{}}\textbf{In}\\1.55\\1.59\end{tabular}} & \makebox[0.9cm]{\begin{tabular}{@{}c@{}}\textbf{Sn}\\1.62\\1.60\end{tabular}} & \makebox[0.9cm]{\begin{tabular}{@{}c@{}}\textbf{Sb}\\1.63\\1.61\end{tabular}} & \makebox[0.9cm]{\textbf{Te}} & \makebox[0.9cm]{\textbf{I}} & \makebox[0.9cm]{\textbf{Xe}} \\
\makebox[0.9cm]{\begin{tabular}{@{}c@{}}\textbf{Cs}\\2.14\\2.26\end{tabular}} & \makebox[0.9cm]{\begin{tabular}{@{}c@{}}\textbf{Ba}\\2.08\\2.03\end{tabular}} & \makebox[0.9cm]{\begin{tabular}{@{}c@{}}\textbf{La}\\1.87\\1.83\end{tabular}} & \makebox[0.9cm]{\begin{tabular}{@{}c@{}}\textbf{Hf}\\1.54\\1.50\end{tabular}} & \makebox[0.9cm]{\begin{tabular}{@{}c@{}}\textbf{Ta}\\1.49\\1.40\end{tabular}} & \makebox[0.9cm]{\begin{tabular}{@{}c@{}}\textbf{W}\\1.43\\1.33\end{tabular}} & \makebox[0.9cm]{\begin{tabular}{@{}c@{}}\textbf{Re}\\1.29\\1.29\end{tabular}} & \makebox[0.9cm]{\begin{tabular}{@{}c@{}}\textbf{Os}\\1.26\\1.27\end{tabular}} & \makebox[0.9cm]{\begin{tabular}{@{}c@{}}\textbf{Ir}\\1.26\\1.29\end{tabular}} & \makebox[0.9cm]{\begin{tabular}{@{}c@{}}\textbf{Pt}\\1.33\\1.33\end{tabular}} & \makebox[0.9cm]{\begin{tabular}{@{}c@{}}\textbf{Au}\\1.40\\1.43\end{tabular}} & \makebox[0.9cm]{\begin{tabular}{@{}c@{}}\textbf{Hg}\\1.52\\1.55\end{tabular}} & \makebox[0.9cm]{\begin{tabular}{@{}c@{}}\textbf{Tl}\\1.63\\1.63\end{tabular}} & \makebox[0.9cm]{\begin{tabular}{@{}c@{}}\textbf{Pb}\\1.73\\1.68\end{tabular}} & \makebox[0.9cm]{\begin{tabular}{@{}c@{}}\textbf{Bi}\\1.71\\1.69\end{tabular}} & \makebox[0.9cm]{\textbf{Po}} & \makebox[0.9cm]{\textbf{At}} & \makebox[0.9cm]{\textbf{Rn}} \\
\makebox[0.9cm]{\textbf{Fr}} & \makebox[0.9cm]{\textbf{Ra}} & \makebox[0.9cm]{\begin{tabular}{@{}c@{}}\textbf{Ac}\\1.97\\1.91\end{tabular}} &  &  &  &  &  &  &  &  &  &  &  &  &  &  &  \\
\midrule
\multicolumn{18}{l}{\textit{Lanthanides:}} \\
 &  &  & \makebox[0.9cm]{\begin{tabular}{@{}c@{}}\textbf{Ce}\\1.79\\1.82\end{tabular}} & \makebox[0.9cm]{\begin{tabular}{@{}c@{}}\textbf{Pr}\\1.83\\1.81\end{tabular}} & \makebox[0.9cm]{\begin{tabular}{@{}c@{}}\textbf{Nd}\\1.81\\1.79\end{tabular}} & \makebox[0.9cm]{\begin{tabular}{@{}c@{}}\textbf{Pm}\\1.81\\1.77\end{tabular}} & \makebox[0.9cm]{\begin{tabular}{@{}c@{}}\textbf{Sm}\\1.76\\1.76\end{tabular}} & \makebox[0.9cm]{\begin{tabular}{@{}c@{}}\textbf{Eu}\\1.90\\1.82\end{tabular}} & \makebox[0.9cm]{\begin{tabular}{@{}c@{}}\textbf{Gd}\\1.74\\1.73\end{tabular}} & \makebox[0.9cm]{\begin{tabular}{@{}c@{}}\textbf{Tb}\\1.72\\1.72\end{tabular}} & \makebox[0.9cm]{\begin{tabular}{@{}c@{}}\textbf{Dy}\\1.70\\1.71\end{tabular}} & \makebox[0.9cm]{\begin{tabular}{@{}c@{}}\textbf{Ho}\\1.68\\1.71\end{tabular}} & \makebox[0.9cm]{\begin{tabular}{@{}c@{}}\textbf{Er}\\1.67\\1.69\end{tabular}} & \makebox[0.9cm]{\begin{tabular}{@{}c@{}}\textbf{Tm}\\1.65\\1.68\end{tabular}} & \makebox[0.9cm]{\begin{tabular}{@{}c@{}}\textbf{Yb}\\1.63\\1.78\end{tabular}} & \makebox[0.9cm]{\begin{tabular}{@{}c@{}}\textbf{Lu}\\1.63\\1.66\end{tabular}} &  \\
\multicolumn{18}{l}{\textit{Actinides:}} \\
 &  &  & \makebox[0.9cm]{\begin{tabular}{@{}c@{}}\textbf{Th}\\1.73\\1.77\end{tabular}} & \makebox[0.9cm]{\begin{tabular}{@{}c@{}}\textbf{Pa}\\--\\1.63\end{tabular}} & \makebox[0.9cm]{\begin{tabular}{@{}c@{}}\textbf{U}\\1.62\\1.60\end{tabular}} & \makebox[0.9cm]{\begin{tabular}{@{}c@{}}\textbf{Np}\\1.62\\1.60\end{tabular}} & \makebox[0.9cm]{\begin{tabular}{@{}c@{}}\textbf{Pu}\\1.63\\1.63\end{tabular}} &  &  &  &  &  &  &  &  &  &  \\
\bottomrule
\end{tabular}
\end{table*}

  \begin{table*}[htbp]
\centering
\scriptsize
\caption{Optimized effective atomic radii (\AA) for L1$_2$ structure.
Upper values: MP (848 compounds, 72 elements),
lower values: OQMD (238 compounds, 58 elements).
Non-metallic elements (H, He, C, N, O, F, Ne, P, S, Cl, Ar, Se, Br, Kr, Te, I, Xe, At, Rn) were excluded from the dataset prior to optimization; only elements appearing in the filtered compound set have numerical values.}
\label{tab:radii_l12}
\setlength{\tabcolsep}{1.5pt}
\begin{tabular}{cccccccccccccccccc}
\toprule
\makebox[0.9cm]{\textbf{H}} &  &  &  &  &  &  &  &  &  &  &  &  &  &  &  &  & \makebox[0.9cm]{\textbf{He}} \\
\midrule
\makebox[0.9cm]{\begin{tabular}{@{}c@{}}\textbf{Li}\\1.43\\1.44\end{tabular}} & \makebox[0.9cm]{\begin{tabular}{@{}c@{}}\textbf{Be}\\1.05\\--\end{tabular}} &  &  &  &  &  &  &  &  &  &  & \makebox[0.9cm]{\begin{tabular}{@{}c@{}}\textbf{B}\\0.50\\--\end{tabular}} & \makebox[0.9cm]{\textbf{C}} & \makebox[0.9cm]{\textbf{N}} & \makebox[0.9cm]{\textbf{O}} & \makebox[0.9cm]{\textbf{F}} & \makebox[0.9cm]{\textbf{Ne}} \\
\makebox[0.9cm]{\begin{tabular}{@{}c@{}}\textbf{Na}\\1.73\\1.77\end{tabular}} & \makebox[0.9cm]{\begin{tabular}{@{}c@{}}\textbf{Mg}\\1.53\\1.49\end{tabular}} &  &  &  &  &  &  &  &  &  &  & \makebox[0.9cm]{\begin{tabular}{@{}c@{}}\textbf{Al}\\1.39\\1.40\end{tabular}} & \makebox[0.9cm]{\begin{tabular}{@{}c@{}}\textbf{Si}\\1.31\\1.28\end{tabular}} & \makebox[0.9cm]{\textbf{P}} & \makebox[0.9cm]{\textbf{S}} & \makebox[0.9cm]{\textbf{Cl}} & \makebox[0.9cm]{\textbf{Ar}} \\
\makebox[0.9cm]{\begin{tabular}{@{}c@{}}\textbf{K}\\2.10\\--\end{tabular}} & \makebox[0.9cm]{\begin{tabular}{@{}c@{}}\textbf{Ca}\\1.87\\1.77\end{tabular}} & \makebox[0.9cm]{\begin{tabular}{@{}c@{}}\textbf{Sc}\\1.54\\1.50\end{tabular}} & \makebox[0.9cm]{\begin{tabular}{@{}c@{}}\textbf{Ti}\\1.42\\1.41\end{tabular}} & \makebox[0.9cm]{\begin{tabular}{@{}c@{}}\textbf{V}\\1.35\\1.36\end{tabular}} & \makebox[0.9cm]{\begin{tabular}{@{}c@{}}\textbf{Cr}\\1.26\\1.41\end{tabular}} & \makebox[0.9cm]{\begin{tabular}{@{}c@{}}\textbf{Mn}\\1.24\\1.27\end{tabular}} & \makebox[0.9cm]{\begin{tabular}{@{}c@{}}\textbf{Fe}\\1.23\\1.11\end{tabular}} & \makebox[0.9cm]{\begin{tabular}{@{}c@{}}\textbf{Co}\\1.18\\1.11\end{tabular}} & \makebox[0.9cm]{\begin{tabular}{@{}c@{}}\textbf{Ni}\\1.08\\1.17\end{tabular}} & \makebox[0.9cm]{\begin{tabular}{@{}c@{}}\textbf{Cu}\\1.18\\1.26\end{tabular}} & \makebox[0.9cm]{\begin{tabular}{@{}c@{}}\textbf{Zn}\\1.34\\1.36\end{tabular}} & \makebox[0.9cm]{\begin{tabular}{@{}c@{}}\textbf{Ga}\\1.43\\1.41\end{tabular}} & \makebox[0.9cm]{\begin{tabular}{@{}c@{}}\textbf{Ge}\\1.45\\1.36\end{tabular}} & \makebox[0.9cm]{\begin{tabular}{@{}c@{}}\textbf{As}\\1.43\\--\end{tabular}} & \makebox[0.9cm]{\textbf{Se}} & \makebox[0.9cm]{\textbf{Br}} & \makebox[0.9cm]{\textbf{Kr}} \\
\makebox[0.9cm]{\begin{tabular}{@{}c@{}}\textbf{Rb}\\2.24\\--\end{tabular}} & \makebox[0.9cm]{\begin{tabular}{@{}c@{}}\textbf{Sr}\\1.97\\2.30\end{tabular}} & \makebox[0.9cm]{\begin{tabular}{@{}c@{}}\textbf{Y}\\1.73\\1.60\end{tabular}} & \makebox[0.9cm]{\begin{tabular}{@{}c@{}}\textbf{Zr}\\1.55\\1.52\end{tabular}} & \makebox[0.9cm]{\begin{tabular}{@{}c@{}}\textbf{Nb}\\1.45\\1.45\end{tabular}} & \makebox[0.9cm]{\begin{tabular}{@{}c@{}}\textbf{Mo}\\1.40\\--\end{tabular}} & \makebox[0.9cm]{\begin{tabular}{@{}c@{}}\textbf{Tc}\\1.35\\--\end{tabular}} & \makebox[0.9cm]{\begin{tabular}{@{}c@{}}\textbf{Ru}\\1.23\\1.27\end{tabular}} & \makebox[0.9cm]{\begin{tabular}{@{}c@{}}\textbf{Rh}\\1.28\\1.26\end{tabular}} & \makebox[0.9cm]{\begin{tabular}{@{}c@{}}\textbf{Pd}\\1.27\\1.29\end{tabular}} & \makebox[0.9cm]{\begin{tabular}{@{}c@{}}\textbf{Ag}\\1.43\\1.43\end{tabular}} & \makebox[0.9cm]{\begin{tabular}{@{}c@{}}\textbf{Cd}\\1.51\\1.50\end{tabular}} & \makebox[0.9cm]{\begin{tabular}{@{}c@{}}\textbf{In}\\1.60\\1.62\end{tabular}} & \makebox[0.9cm]{\begin{tabular}{@{}c@{}}\textbf{Sn}\\1.62\\1.64\end{tabular}} & \makebox[0.9cm]{\begin{tabular}{@{}c@{}}\textbf{Sb}\\1.64\\1.42\end{tabular}} & \makebox[0.9cm]{\textbf{Te}} & \makebox[0.9cm]{\textbf{I}} & \makebox[0.9cm]{\textbf{Xe}} \\
\makebox[0.9cm]{\begin{tabular}{@{}c@{}}\textbf{Cs}\\2.26\\--\end{tabular}} & \makebox[0.9cm]{\begin{tabular}{@{}c@{}}\textbf{Ba}\\2.12\\--\end{tabular}} & \makebox[0.9cm]{\begin{tabular}{@{}c@{}}\textbf{La}\\1.78\\1.75\end{tabular}} & \makebox[0.9cm]{\begin{tabular}{@{}c@{}}\textbf{Hf}\\1.52\\1.51\end{tabular}} & \makebox[0.9cm]{\begin{tabular}{@{}c@{}}\textbf{Ta}\\1.43\\1.47\end{tabular}} & \makebox[0.9cm]{\begin{tabular}{@{}c@{}}\textbf{W}\\1.37\\--\end{tabular}} & \makebox[0.9cm]{\begin{tabular}{@{}c@{}}\textbf{Re}\\1.37\\--\end{tabular}} & \makebox[0.9cm]{\begin{tabular}{@{}c@{}}\textbf{Os}\\1.30\\--\end{tabular}} & \makebox[0.9cm]{\begin{tabular}{@{}c@{}}\textbf{Ir}\\1.31\\1.30\end{tabular}} & \makebox[0.9cm]{\begin{tabular}{@{}c@{}}\textbf{Pt}\\1.31\\1.33\end{tabular}} & \makebox[0.9cm]{\begin{tabular}{@{}c@{}}\textbf{Au}\\1.43\\1.43\end{tabular}} & \makebox[0.9cm]{\begin{tabular}{@{}c@{}}\textbf{Hg}\\1.57\\1.57\end{tabular}} & \makebox[0.9cm]{\begin{tabular}{@{}c@{}}\textbf{Tl}\\1.67\\1.67\end{tabular}} & \makebox[0.9cm]{\begin{tabular}{@{}c@{}}\textbf{Pb}\\1.68\\1.72\end{tabular}} & \makebox[0.9cm]{\begin{tabular}{@{}c@{}}\textbf{Bi}\\1.72\\1.27\end{tabular}} & \makebox[0.9cm]{\textbf{Po}} & \makebox[0.9cm]{\textbf{At}} & \makebox[0.9cm]{\textbf{Rn}} \\
\makebox[0.9cm]{\textbf{Fr}} & \makebox[0.9cm]{\textbf{Ra}} & \makebox[0.9cm]{\begin{tabular}{@{}c@{}}\textbf{Ac}\\1.91\\--\end{tabular}} &  &  &  &  &  &  &  &  &  &  &  &  &  &  &  \\
\midrule
\multicolumn{18}{l}{\textit{Lanthanides:}} \\
 &  &  & \makebox[0.9cm]{\begin{tabular}{@{}c@{}}\textbf{Ce}\\1.69\\1.74\end{tabular}} & \makebox[0.9cm]{\begin{tabular}{@{}c@{}}\textbf{Pr}\\1.79\\1.71\end{tabular}} & \makebox[0.9cm]{\begin{tabular}{@{}c@{}}\textbf{Nd}\\1.79\\1.66\end{tabular}} & \makebox[0.9cm]{\begin{tabular}{@{}c@{}}\textbf{Pm}\\1.74\\--\end{tabular}} & \makebox[0.9cm]{\begin{tabular}{@{}c@{}}\textbf{Sm}\\1.76\\1.69\end{tabular}} & \makebox[0.9cm]{\begin{tabular}{@{}c@{}}\textbf{Eu}\\1.85\\1.68\end{tabular}} & \makebox[0.9cm]{\begin{tabular}{@{}c@{}}\textbf{Gd}\\1.75\\1.59\end{tabular}} & \makebox[0.9cm]{\begin{tabular}{@{}c@{}}\textbf{Tb}\\1.72\\1.65\end{tabular}} & \makebox[0.9cm]{\begin{tabular}{@{}c@{}}\textbf{Dy}\\1.72\\1.59\end{tabular}} & \makebox[0.9cm]{\begin{tabular}{@{}c@{}}\textbf{Ho}\\1.72\\1.59\end{tabular}} & \makebox[0.9cm]{\begin{tabular}{@{}c@{}}\textbf{Er}\\1.68\\1.58\end{tabular}} & \makebox[0.9cm]{\begin{tabular}{@{}c@{}}\textbf{Tm}\\1.67\\1.58\end{tabular}} & \makebox[0.9cm]{\begin{tabular}{@{}c@{}}\textbf{Yb}\\1.65\\1.58\end{tabular}} & \makebox[0.9cm]{\begin{tabular}{@{}c@{}}\textbf{Lu}\\1.64\\1.56\end{tabular}} &  \\
\multicolumn{18}{l}{\textit{Actinides:}} \\
 &  &  & \makebox[0.9cm]{\begin{tabular}{@{}c@{}}\textbf{Th}\\1.76\\1.71\end{tabular}} & \makebox[0.9cm]{\begin{tabular}{@{}c@{}}\textbf{Pa}\\1.58\\1.60\end{tabular}} & \makebox[0.9cm]{\begin{tabular}{@{}c@{}}\textbf{U}\\1.49\\1.58\end{tabular}} & \makebox[0.9cm]{\begin{tabular}{@{}c@{}}\textbf{Np}\\1.53\\1.61\end{tabular}} & \makebox[0.9cm]{\begin{tabular}{@{}c@{}}\textbf{Pu}\\1.56\\1.61\end{tabular}} &  &  &  &  &  &  &  &  &  &  \\
\bottomrule
\end{tabular}
\end{table*}



\subsection{Analysis of B2-L1$_2$ Radius Difference}

本研究が扱う重要な問題は、B2(CN=8)構造とL1$_2$(CN=12)構造の間で原子半径がどれだけ変化するかである。Goldschmidtの経験則によれば、配位数が12から8に変化すると、半径は約3％縮小すると考えられる。この予測は、同一元素について測定されたB2とL1$_2$の半径を比較することで検証できる。

表~\ref{tab:top_contraction}は、MPデータセットにおいてB2-L1$_2$半径差(収縮率)が最も大きい元素を列挙している。最大の収縮を示すのはRe($-5.9$\%)であり、次いでCa($-5.6$\%), Hg($-3.6$\%)が続く。一方、Ni, Pu, Uなどの元素はB2半径がL1$_2$半径より大きい逆の傾向を示している。

% A key question addressed by this study is how effective atomic radii change between B2 (CN=8) and L1$_2$ (CN=12) structures. According to Goldschmidt's empirical rule, a coordination number change from 12 to 8 should result in approximately 3\% radius contraction.

% Table~\ref{tab:top_contraction} lists the elements showing the largest B2-L1$_2$ radius differences (contraction rates) in the MP dataset. The largest contraction is observed for B (boron, $-8.7$\%), followed by Li ($-7.5$\%). Transition metals such as Re and Tc also show contractions of 5--6\%. Conversely, elements such as Ni, U, and Cu show the opposite trend, with B2 radii larger than L1$_2$ radii.

\begin{table*}[htbp]
\small
    \centering
    \caption{Elements with the largest B2-L1$_2$ radius differences in the MP dataset. Negative contraction indicates B2 radius smaller than L1$_2$ radius. The top five contracting and top five expanding elements are shown.}
    \label{tab:top_contraction}
\scalebox{0.9}{
    \begin{tabular}{lccccc}
        \toprule
        Element & B2 radius (\AA) & L1$_2$ radius (\AA) & Diff. (\AA) & Contraction (\%) & Group \\
        \midrule
        Re & 1.29 & 1.37 & $-$0.08 & $-$5.9 & 5d transition \\
        Ca & 1.75 & 1.86 & $-$0.10 & $-$5.6 & Alkaline earth \\
        Hg & 1.52 & 1.57 & $-$0.06 & $-$3.6 & 5d transition \\
        Ir & 1.27 & 1.31 & $-$0.04 & $-$3.4 & 5d transition \\
        Ga & 1.35 & 1.40 & $-$0.05 & $-$3.4 & Post-transition \\
        \midrule
        Ce & 1.79 & 1.69 & $+$0.10 & $+$6.0 & Lanthanide \\
        Mn & 1.31 & 1.23 & $+$0.08 & $+$6.4 & 3d transition \\
        U & 1.62 & 1.52 & $+$0.10 & $+$6.5 & Actinide \\
        Pu & 1.63 & 1.52 & $+$0.11 & $+$7.4 & Actinide \\
        Ni & 1.18 & 1.09 & $+$0.09 & $+$7.9 & 3d transition \\
        \bottomrule
    \end{tabular}
    }
\end{table*}

図~\ref{fig:B2_L12_scatter}は、両方の構造タイプにデータを持つ全68元素のB2半径とL1$_2$半径の散布図を示す。元素は対角線の両側に分布しており、B2半径がL1$_2$半径より小さい元素(B, Li, Re等)と大きい元素(Ni, U, Cu等)の両方が見られる。

% A scatter plot of B2 vs. L1$_2$ radii for all 68 elements with data in both structure types is shown in Fig.~\ref{fig:B2_L12_scatter}. Elements are distributed on both sides of the diagonal, with some having smaller B2 radii (B, Li, Re) and others having larger B2 radii (Ni, U, Cu).

\begin{figure}[htbp]
    \centering
    \includegraphics[width=0.4\textwidth]{07_B2_vs_L12_scatter.png}
    \caption{Scatter plot of B2 vs. L1$_2$ effective radii (MP) for 68 elements present in both datasets. The dashed line represents equal radii. Elements are distributed on both sides, with B and Li showing contraction in B2, while Ni and U show expansion.}
    \label{fig:B2_L12_scatter}
\end{figure}

全元素におけるB2-L1$_2$収縮率の分布を図~\ref{fig:B2_contraction}に示す。Goldschmidtの3\%ルール(赤破線で示される)との比較では、max基準による格子定数モデルの採用により、従来の平均化モデルと比べて収縮率の分布が異なる傾向を示す。
% The distribution of B2-L1$_2$ contraction rates across all elements is shown in Fig.~\ref{fig:B2_contraction}. The majority of transition metals show contractions in the range of 0--5\%, reasonably consistent with Goldschmidt's 3\% rule (indicated by the red dashed line). The distribution is asymmetric, with a longer negative tail (B2 $<$ L1$_2$) than positive tail (B2 $>$ L1$_2$), indicating that elements systematically contract in B2 more often than they expand.

\begin{figure}[htbp]
    \centering
    \includegraphics[width=0.4\textwidth]{gs_04_B2_contraction.png}
    \caption{Distribution of B2-L1$_2$ radius contraction rates for all elements (MP data). The red dashed line indicates Goldschmidt's predicted 3\% contraction for CN 12 $\rightarrow$ 8. Most transition metals fall near this value, while some elements show anomalously large contractions.}
    \label{fig:B2_contraction}
\end{figure}

\subsection{Element Group Analysis}
有効半径の体系的傾向をみるため、元素をアルカリ土類、3d/4d/5d遷移金属、ランタノイド、アクチノイド、遷移後金属に分類してB2--L1$_2$差を比較した（アルカリ金属は本研究の解析対象から除外済み）。表~\ref{tab:element_groups}に示すように、3d遷移金属の平均変化率は$+1.7$\%、4d遷移金属は$+0.7$\%、5d遷移金属は$-0.3$\%であり、max基準モデルではB2半径がL1$_2$半径より大きくなる元素が多く見られる。ランタノイドは平均$+0.5$\%(14元素)、アクチノイドは$+4.1$\%とB2の方が大きい傾向が見られた。遷移後金属(Al, Ga, In等)は平均$-1.9$\%の収縮を示した。

\begin{table*}[htbp]      \small
    \centering
    \caption{B2-L1$_2$ radius contraction statistics by element group (MP data). The ``Avg.'' column shows the mean contraction rate, and ``Range'' shows the minimum to maximum values within each group. Negative values indicate B2 radius smaller than L1$_2$ radius.}
    \label{tab:element_groups}
\scalebox{0.9}{
    \begin{tabular}{lcccccc}
        \toprule
        Element group & N & B2 avg. (\AA) & L1$_2$ avg. (\AA) & Avg. (\%) & Range (\%) & Characteristic \\
        \midrule
        Alkaline earth & 5 & 1.68 & 1.69 & $-$0.2 & $-$5.6 to $+$3.7 & Ca contraction \\
        3d transition & 10 & 1.30 & 1.28 & $+$1.7 & $-$2.8 to $+$7.9 & Ni expand \\
        4d transition & 10 & 1.43 & 1.42 & $+$0.7 & $-$2.0 to $+$4.0 & Mixed \\
        5d transition & 10 & 1.44 & 1.44 & $-$0.3 & $-$5.9 to $+$5.3 & Re contraction \\
        Lanthanides & 14 & 1.73 & 1.73 & $+$0.5 & $-$2.2 to $+$6.0 & Slight expansion \\
        Actinides & 5 & 1.72 & 1.65 & $+$4.1 & $-$1.3 to $+$7.4 & B2 expansion \\
        Post-transition & 6 & 1.54 & 1.57 & $-$1.9 & $-$3.4 to $-$0.5 & Moderate contraction \\
        \bottomrule
    \end{tabular}
    }
\end{table*}

\subsection{Database Consistency}
両データベースに共通して存在する化合物について比較すると、B2では357件、L1$_2$では186件の共通化合物が存在した。格子定数の比較では、B2で$R^2 = 0.986$、L1$_2$で$R^2 = 0.986$と高い整合性が認められた(表~\ref{tab:consistency})。有効半径についても、B2では68の共通元素に対して$R^2 = 0.969$と高い整合性が確認された。L1$_2$の半径整合性は$R^2 = 0.885$(58共通元素)であり、B2よりやや低いが、max基準モデルの採用により従来の平均化モデル($R^2 = 0.657$)から改善された。

% Comparing compounds present in both databases, 357 common B2 compounds and 186 common L1$_2$ compounds were identified. Lattice constant comparisons showed high consistency with $R^2 = 0.986$ for B2 and $R^2 = 0.986$ for L1$_2$ (Table~\ref{tab:consistency}). For effective radii, B2 showed high consistency ($R^2 = 0.969$) across 68 common elements. L1$_2$ radius consistency was $R^2 = 0.885$ across 58 common elements, improved from the previous averaging model ($R^2 = 0.657$) by adopting the max-based model.

\begin{table}[htbp]      \small
    \centering
    \caption{Consistency between MP and OQMD results for compounds present in both databases. The lattice constant $R^2$ compares DFT-calculated values between databases, while the radius $R^2$ compares the optimized effective radii for common elements.}
    \label{tab:consistency}
\scalebox{0.9}{
    \begin{tabular}{lcccc}
        \toprule
        Structure & \shortstack{Common\\compounds} & \shortstack{Lattice\\$R^2$} & \shortstack{Common\\elements} & \shortstack{Radius\\$R^2$} \\
        \midrule
        B2 & 357 & 0.986 & 68 & 0.969 \\
        L1$_2$ & 186 & 0.986 & 58 & 0.885 \\
        \bottomrule
    \end{tabular}
    }
\end{table}

\begin{figure*}[t]
  \centering
  \begin{minipage}[t]{0.48\textwidth}
    \centering
    \includegraphics[width=0.9\linewidth]{a7_common_B2.png}
    \caption{Comparison of lattice constants between MP and OQMD for 357 B2 compounds present in both databases ($R^2$ = 0.986).}
    \label{fig:common_B2}
  \end{minipage}
  \hfill
  \begin{minipage}[t]{0.48\textwidth}
    \centering
    \includegraphics[width=0.9\linewidth]{a7_common_L12.png}
    \caption{Comparison of lattice constants between MP and OQMD for 186 L1$_2$ compounds present in both databases ($R^2$ = 0.986).}
    \label{fig:common_L12}
  \end{minipage}
\end{figure*}


%%%%%%%%%%%%%%%%%%%%%%%%%%%%%%%%%%%%%
\section{Discussion}

\subsection{Effect of Non-Metallic and Alkali Metal Exclusion}

非金属元素およびアルカリ金属元素(Li, Na, K, Rb, Cs)を含む化合物の除外は硬球モデルの精度に影響する。非金属元素はイオン性・共有結合性が強く、硬球近似が不適切である。アルカリ金属は大きなイオン半径と強いイオン性のため、B2構造において格子定数の系統的な過大評価を引き起こす。表~\ref{tab:nonmetal_effect}に示すように、改善効果はケースによって異なる。なお、Before RMSEは従来のA--B接触のみモデル（非金属・アルカリ金属を含む）による値であり、After RMSEは本研究のmax基準反復最適化モデル（非金属・アルカリ金属除外後）による値であるため、改善率にはモデル変更とデータフィルタリングの両方の効果が含まれる。

% The most important methodological finding is the pronounced effect of excluding non-metallic elements on hard-sphere model accuracy. As shown in Table~\ref{tab:nonmetal_effect}, the RMSE improvement varies significantly across cases. MP-B2 shows a dramatic 67\% improvement (0.158~\AA\ to 0.052~\AA), demonstrating that strongly covalent/ionic compounds such as IrN, OsN, TiO, and CuO deviated substantially from the hard-sphere model.

\begin{table}[htbp]
\small
    \centering
    \caption{Effect of non-metallic and alkali metal exclusion on lattice constant prediction accuracy.}
    \label{tab:nonmetal_effect}
\scalebox{0.9}{
    \begin{tabular}{lccc}
        \toprule
        Case & Before RMSE (\AA) & After RMSE (\AA) & Improvement \\
        \midrule
        MP --- B2 & 0.158 & 0.054 & 66\% \\
        MP --- L1$_2$ & 0.206 & 0.116 & 44\% \\
        OQMD --- B2 & 0.102 & 0.095 & 7\% \\
        OQMD --- L1$_2$ & 0.067 & 0.066 & 1\% \\
        \bottomrule
    \end{tabular}
    }
\end{table}

この結果は、硬球幾何モデルの適用範囲を明確に示している。非金属元素・アルカリ金属元素を除いた金属間化合物に対しては良好な予測を提供するが、これらの元素を含む化合物では電荷移動、共有結合性、強いイオン性により原子間距離が硬球モデルから逸脱する。

\subsection{Structure Dependence of Effective Radii}

B2--L1$_2$半径差の元素群分析(表~\ref{tab:element_groups})から、以下の系統的傾向が明らかになった。

元素群別の平均B2--L1$_2$半径差は、3d遷移金属で$+1.7$\%(10元素)、4d遷移金属で$+0.7$\%(10元素)、5d遷移金属で$-0.3$\%(10元素)であった。これらのB2-L1$_2$差は正負両方に分布しており、Goldschmidtの一律3\%収縮則とは合致しない。個々の元素レベルでは、Re($-5.9$\%)やCa($-5.6$\%)のように大きな収縮を示す元素がある一方、Ni($+7.9$\%)やPu($+7.4$\%)のようにB2半径がL1$_2$半径より大きい元素も存在する。これは配位数変化のみでは説明できない元素固有の効果を示している。

ランタノイドの平均収縮率は$+0.2$\%(13元素)であった。4f電子の遮蔽効果によりランタノイドでは配位環境依存性が小さい。なお、アルカリ金属元素(Li, Na, K, Rb, Cs)は強いイオン性のため硬球モデルへの適合が悪く、本研究では解析対象から除外した。

% Element-group average B2--L1$_2$ radius differences: 3d: $+2.1$\% (10 elements), 4d: $+0.2$\% (10 elements), 5d: $-0.4$\% (10 elements). Lanthanides: $+0.8$\% (13 elements), alkali metals: $-1.1$\% (3 elements).

\subsection{Implications for Hume-Rothery Rules}
Hume--Rothery則のサイズ因子は、通常、元素ごとに単一の原子半径を仮定して適用される~\citep{humerothery1936, miracle2017}。しかし本研究の結果は、B2とL1$_2$のような秩序相では、有効半径が構造依存であることを示している。遷移金属の大半ではB2-L1$_2$差は数\%以内であり実用上の影響は限定的であるが、特定の元素(Ni: $+7.9$\%, Pu: $+7.4$\%, Re: $-5.9$\%等)では約8\%に達する差が生じるため、これらの元素を含む系では構造に応じた半径の使い分けが重要となる。

したがって、相安定性評価や合金設計では、想定する構造に応じた半径を使い分ける方が妥当である。L1$_2$についてはGoldschmidt半径が参考になる($R^2 = 0.858$)一方、B2については本研究で決定した構造特異的半径を用いる意義が大きい。

この知見は、Hume--Rothery型のサイズ因子による議論だけでなく、機械学習を用いた材料特性予測において原子半径を特徴量として利用する場合にも重要である。とくに、結晶構造や局所結合環境に応じて有効半径が大きく変化する元素では、単一の原子半径を記述子として用いると、環境依存的な効果を見落とす可能性がある。そのような系では、構造特異的、あるいは局所環境を反映した記述子を用いることが望ましい。


\subsection{Comparison with Previous Studies}

いくつかの先行研究では、金属間化合物のデータから有効原子半径を決定しようと試みてきた。しかし、現代の第一原理データベースから得られるような大規模なデータセットを用いた研究は存在しない。

King~~\citep{king1966} は、Laves相やその他の金属間構造の格子定数を分析し、有効半径を導出した。その結果、Goldschmidt則と同様の配位数依存性の証拠を見出した。しかし、彼のデータセットは実験的に特徴づけられた化合物に限定され、数千ではなく数百程度であった。Teatum et al.~~\citep{teatum1960} は、実験データと理論的考察を組み合わせて、金属間化合物の構造予測に用いる原子半径をまとめた。彼らの値は広く用いられてきたが、構造固有の効果を考慮していない。近年では機械学習手法が金属間化合物の格子定数やその他の特性の予測に応用されている ~\citep{ward2016}。しかしこれらは通常、原子半径を最適化可能なパラメータではなく固定入力特徴として扱う。

本研究では、大規模な第一原理データ(総計3511件)を用いて構造特異的有効半径を決定し、配位数効果や元素群の傾向を分析した。

% Several previous studies have attempted to determine effective atomic radii from intermetallic compound data, though none have used datasets as large as those available from modern first-principles databases.

% Using large-scale first-principles data (3982 compounds in total), structure-specific effective radii were determined, enabling analysis of coordination number effects and element-group trends.

\subsection{Limitations and Sources of Uncertainty}

Several limitations should be considered when interpreting our results:

\textbf{DFT systematic errors}: 
% MPとOQMDの両方の計算で用いられるGGA-PBE関数は、実験値と比較して格子定数を1～2\%系統的に過大評価することが知られている。
The GGA-PBE functional used in both MP and OQMD calculations is known to systematically overestimate lattice constants by 1--2\% compared to experiment. This systematic error affects all compounds similarly and thus should not significantly impact the relative radii determined in this study, but the absolute scale of our radii may be slightly larger than would be obtained from experimental data.

\textbf{温度効果}: DFT計算は0~K(零点運動を無視)に対応するが、実験測定や古典半径の表は通常室温を指す。熱膨張効果は一般的に小さい(数パーセント)が、我々の半径と古典値の差異に寄与する可能性がある。

\textbf{磁気効果}: 磁性元素(Fe、Co、Ni、Mn、Cr、希土類元素)では、磁気状態が原子体積に影響し、結果として有効半径を変える。MP法やOQMD法ではスピン分極を考慮するが、磁気秩序の扱いは必ずしも基底状態に対応しない。特に反強磁性や複雑な磁気構造の場合に顕著である。

\textbf{モデルの限界}: 剛体球幾何モデルは電荷移動、混成軌道、dバンド充填などの電子効果を無視している。非金属元素やアルカリ金属を含む化合物での大きな予測誤差は、これらの電子効果が無視できないことを示しており、本研究ではこれらの元素を解析対象から除外した。

\textbf{データの品質}: 非金属元素・アルカリ金属元素除外と格子定数範囲フィルタを適用したが、データベース内の一部のエントリには構造決定や計算パラメータの誤りが含まれる可能性がある。B2ではMPとOQMDの半径に高い一貫性($R^2 = 0.955$)があり、誤りは系統的ではないと考えられる。L1$_2$の半径整合性($R^2 = 0.854$)はB2より低く、OQMDのL1$_2$データ数(235件)がMP(723件)より少ないことが一因である。データポイントが少ない元素では半径の不確実性が大きくなる。

\textbf{max基準モデルの近似}: 本研究ではB2(式~\ref{eq:B2})およびL1$_2$(式~\ref{eq:L12})の両方にmax基準モデルを採用した。これは剛体球が重ならない最小格子定数を与える物理的に明確なモデルである。しかし、実際の結晶では原子は完全な剛体球ではなく、電子的効果による緩和が存在する。B2のmax基準モデルでは、同種原子間接触(A--AまたはB--B)が支配的な化合物において、非支配的な原子種の半径への制約が弱まる。この問題は不等式制約付きSLSQP最適化により緩和され、非支配的接触条件が観測格子定数以下となることが保証される。

% \textbf{Temperature effects}: DFT calculations correspond to 0~K (with zero-point motion neglected), while experimental measurements and classical radius tabulations typically refer to room temperature.

% \textbf{Magnetic effects}: For magnetic elements (Fe, Co, Ni, Mn, Cr, and rare earths), the magnetic state can affect atomic volume and thus effective radius.

% \textbf{Model limitations}: The hard-sphere geometric models used here neglect electronic effects such as charge transfer, hybridization, and d-band filling that can influence interatomic distances.

% \textbf{Data quality}: While we applied non-metallic element exclusion and lattice constant range filters, some entries in the databases may have errors in structure assignment or computational parameters.

% \textbf{L1$_2$ max-based model approximation}: The max-based model (Eq.~\ref{eq:L12}) gives the minimum lattice parameter that avoids hard-sphere overlap. However, real atoms are not perfect hard spheres, and electronic relaxation effects exist. Thus, the max criterion may not always provide the optimal description.


\subsection{Validation Against Multi-Component Alloys}

本研究で決定した構造特異的有効半径の多成分合金への適用性を検証するため、文献で報告されているHEA（高エントロピー合金）およびBMG（バルク金属ガラス）の実験データとの比較を行った。

\subsubsection{HEA Lattice Constant Prediction}

Vegard則に基づき、BCC型HEAには$a = 4\bar{r}/\sqrt{3}$（$\bar{r} = \sum c_i r_i$）を用いた。FCC型HEAには、有効半径の決定に用いたmax接触モデルとの整合性を保つため、ペアワイズmax接触モデル$a = 2\sqrt{2}\sum_i \sum_j c_i c_j \max(r_i, r_j)$を採用した。12種のHEA実験値と比較した結果を表~\ref{tab:hea_validation}にまとめる。

BCC HEAに対してはB2半径が有効であり、OQMD-B2半径ではRMSE = 0.051~\AA（平均相対誤差1.54\%）を達成した。これはPauling半径（RMSE = 0.115~\AA）およびGoldschmidt半径（RMSE = 0.061~\AA）より高精度である。AlCoCrFeNi（BCC）では誤差0.51\%と特に良好な一致を示した。

FCC HEAに対しては、単純Vegard式（$a = 2\sqrt{2}\bar{r}$）では有効半径が格子定数を系統的に過小評価するが、ペアワイズmax接触モデルの導入により大幅に改善された。MP-B2半径ではRMSE = 0.017~\AA（平均相対誤差0.36\%）を達成し、Pauling半径（RMSE = 0.069~\AA）を凌駕する。この改善は、有効半径フィッティングにおけるL1$_2$化合物の約51\%が異種原子接触（$a = \sqrt{2}(r_A + r_B)$）で格子定数が決定されることと整合的である。

\begin{table}[htbp]
\small
    \centering
    \caption{HEA lattice constant prediction using structure-matched effective radii. BCC HEAs are predicted with $a = 4\bar{r}/\sqrt{3}$, FCC HEAs with the pairwise max-contact model $a = 2\sqrt{2}\sum_i \sum_j c_i c_j \max(r_i, r_j)$. Representative entries are shown.}
    \label{tab:hea_validation}
\scalebox{0.82}{
    \begin{tabular}{llcccl}
        \toprule
        HEA & Structure & $a_{\mathrm{exp}}$ (\AA) & $a_{\mathrm{pred}}$ (\AA) & Rel.\ err.\ (\%) & Radius set \\
        \midrule
        AlCoCrFeNi & BCC & 2.871 & 2.886 & 0.51 & MP-B2 \\
        TiZrHfNbTa & BCC & 3.404 & 3.471 & 1.98 & MP-B2 \\
        NbMoTaW & BCC & 3.213 & 3.158 & 1.70 & OQMD-B2 \\
        VNbMoTaW & BCC & 3.183 & 3.137 & 1.46 & OQMD-B2 \\
        TiZrNbMoTa & BCC & 3.293 & 3.255 & 1.14 & OQMD-B2 \\
        \midrule
        CoCrFeMnNi & FCC & 3.592 & 3.606 & 0.39 & OQMD-L1$_2$ \\
        CoCrFeNi & FCC & 3.572 & 3.568 & 0.10 & OQMD-L1$_2$ \\
        CoCrFeNiCu & FCC & 3.590 & 3.603 & 0.36 & OQMD-L1$_2$ \\
        \bottomrule
    \end{tabular}
    }
\end{table}

\subsubsection{BMG Glass-Forming Ability Indicators}

液体/アモルファス相は局所的にCN $\approx$ 12の配位環境を持つため、BMG合金の原子サイズミスマッチ$\delta$の算出にはL1$_2$半径を使用した。10種のBMG合金および5種の結晶性HEAについて$\delta$を計算した結果、BMG形成合金は$\delta > 6$\%、結晶性HEAは$\delta < 6$\%の領域に明確に分離された。

Inoueの三経験則（$N \geq 3$、大きな原子サイズ差、負の混合エンタルピー）との比較では、Vitreloy 1（Zr$_{41}$Ti$_{14}$Cu$_{13}$Ni$_{10}$Be$_{22}$, $\delta = 16.4$\%）やZr$_{55}$Cu$_{30}$Al$_{10}$Ni$_5$（$\delta = 12.9$\%）など高GFA合金が3条件すべてを満たす一方、ガラス形成能の低いAu$_{77}$Ge$_{14}$Si$_9$は$\delta = 2.5$\%と小さかった。

これらの結果は、二元系化合物から決定した有効半径が、多成分合金のサイズミスマッチ評価にも有効であることを示しており、特にB2半径のBCC HEAへの適用において古典半径を上回る予測精度が得られた。


%%%%%%%%%%%%%%%%%%%%%%%%%%%%%%%%%%%%%
\section{Conclusions}
We determined effective atomic radii for B2 and L1$_2$ binary intermetallic compounds by optimizing against lattice constant data from the Materials Project and OQMD databases. Non-metallic and alkali metal elements were excluded and a lattice constant range filter (2.0--8.0~\AA) was applied to define the applicability domain of the hard-sphere model. A total of 3511 compounds (MP: 346 B2, 723 L1$_2$; OQMD: 2207 B2, 235 L1$_2$) were analyzed using inequality-constrained SLSQP optimization with Bootstrap uncertainty estimation. The main conclusions are as follows.
\begin{enumerate}
    \item \textbf{Predictive accuracy}: the optimized radii reproduce lattice constants with RMSE = 0.054~\AA\ for MP B2, 0.095~\AA\ for OQMD B2, 0.116~\AA\ for MP L1$_2$, and 0.066~\AA\ for OQMD L1$_2$, corresponding to $R^2 = 0.935$--0.975.
    \item \textbf{Effectiveness of element exclusion}: removal of compounds containing non-metallic elements and alkali metals improved prediction accuracy, particularly for MP B2 (RMSE reduced by 66\%) and MP L1$_2$ (44\%), confirming that the hard-sphere model is appropriate for metallic intermetallic compounds but not for covalent/ionic systems or compounds with strong ionic character.
    \item \textbf{Structure dependence of effective radii}: L1$_2$ radii correlate with Goldschmidt metallic radii ($R^2 = 0.858$), whereas B2 radii display element-specific deviations. With the max-based models for both B2 and L1$_2$, the B2-L1$_2$ radius differences are distributed on both sides: elements such as Re ($-5.9$\%) and Ca ($-5.6$\%) show contraction, while Ni ($+7.9$\%) and Pu ($+7.4$\%) show expansion.
    \item \textbf{Database consistency}: MP and OQMD give highly consistent lattice constants for common compounds ($R^2 > 0.985$) and consistent B2 radii ($R^2 = 0.955$), indicating that the optimized radii are stable against reasonable differences in first-principles workflows. L1$_2$ radius consistency ($R^2 = 0.854$) is lower, partly due to smaller OQMD L1$_2$ dataset size.
    \item \textbf{Multi-component alloy validation}: B2 radii predict BCC HEA lattice constants with RMSE = 0.051~\AA\ (OQMD-B2), outperforming Pauling radii (0.115~\AA). With the pairwise max-contact model, FCC HEA lattice constants are predicted with RMSE = 0.017~\AA\ (MP-B2), surpassing even Pauling radii (0.069~\AA). For BMG alloys, the atomic size mismatch $\delta$ calculated from these radii clearly separates glass-forming ($\delta > 6$\%) from crystalline ($\delta < 6$\%) compositions.
\end{enumerate}

The tabulated radii provide a practical reference for lattice-constant estimation, size-mismatch analysis, and phase-stability discussions in B2- and L1$_2$-containing alloy systems.


%%%%%%%%%%%%%%%%%%%%%%%%%%%%%%%%%%%%%
\section*{CRediT authorship contribution statement}

\textbf{Satoshi Minamoto}: Conceptualization, Methodology, Software, Validation, Formal analysis, Investigation, Data curation, Writing - original draft, Writing - review \& editing, Visualization.


%%%%%%%%%%%%%%%%%%%%%%%%%%%%%%%%%%%%%
\section*{Declaration of competing interest}

The authors declare that they have no known competing financial interests or personal relationships that could have appeared to influence the work reported in this paper.

%%%%%%%%%%%%%%%%%%%%%%%%%%%%%%%%%%%%%
\section*{Data availability}

The optimized effective radii and analysis scripts are available at MDR, NIMS Data Repository. Raw data were obtained from the Materials Project (https://materialsproject.org) and OQMD (http://oqmd.org) under their respective data use agreements.


%%%%%%%%%%%%%%%%%%%%%%%%%%%%%%%%%%%%%
\section*{Acknowledgments}

The author acknowledges the Materials Project and OQMD teams for providing open access to their computational materials databases, which made this study possible. 



%%%%%%%%%%%%%%%%%%%%%%%%%%%%%%%%%%%%%
\bibliographystyle{elsarticle-num}
\bibliography{minamoto}




\end{document}
